\subsection{محدودیت‌ها و تحلیل شکست}
علیرغم عملکرد برجسته در بنچ‌مارک‌های استاندارد، محدودیت‌های زیر برای \lr{FLAP} قابل شناسایی است:
\begin{itemize}
    \item \textbf{افت دقت در مدل‌های بزرگ علّی (\lr{Qwen}):} در مدل‌های پارامتربالا مانند \lr{Qwen}، افت دقت به سطح \lr{60.78\%} مشاهده می‌شود که ناشی از حساسیت فاحش دیکودرهای خودبازگشت به تنکی ستونی بدون تنظیم مجدد ماتریس‌های کلید-مقدار در لایه‌های عمیق است.
    \item \textbf{وابستگی به توزیع داده‌های کالیبراسیون:} محاسبه دقیق نوسانات نیازمند نمونه‌های کالیبراسیون با تنوع زبانی بالا است؛ در صورتی که داده‌های کالیبراسیون محدود یا دارای بایاس باشند، واریانس ویژگی‌ها گمراه‌کننده شده و منجر به هرس نامناسب کانال‌های مهم اطلاعاتی می‌گردد.
    \item \textbf{نیاز به کامپایلر سخت‌افزاری جهت بهره‌برداری کامل از تنکی:} هرس ستونی ماسک‌شده، در صورت عدم حذف فیزیکی ابعاد ماتریس یا عدم استفاده از هسته‌های بهینه‌شده \lr{Sparse Tensor Core}، لزوماً به کاهش خطی زمان استنتاج خام در تمامی سخت‌افزارها منتج نخواهد شد.
\end{itemize}
